\section{Discussion}
\label{sec:discussion}

\subsection{The Threshold Is Missed on Both Forms, but for Different Reasons}
\label{subsec:disc-finding}

Read strictly, the answer to the research question is negative on both forms:
GPT-4o mini reaches $\kappa = 0.578$ on Raw and $\kappa = 0.335$ on Improved,
neither meeting the pre-registered bar of $\kappa \geq 0.70$. But the two
shortfalls have different characters. On Raw the gap is narrow and symmetric:
the model sits just below \emph{substantial} agreement and close to the
$\kappa_{\text{control}} = 0.674$ measured between the two human annotators, and
its errors do not lean in any direction. On Improved the gap is wide and
one-sided: agreement is only \emph{fair}, far below the human ceiling on that
form ($0.556$), and driven by a significant optimistic bias. The Raw result is
consistent with a capable-but-imperfect evaluator; the Improved result is
consistent with an evaluator that is being misled.

\subsection{The Central Finding: Agreement Degrades on AI-Improved Reports}
\label{subsec:disc-gap}

The most important result of this study is the \emph{difference} between the two
forms. The same model, prompt, rubric, and scoring function, applied to the same
139 issues, agrees with humans far less well once those issues have been
rewritten by another LLM. The degradation is concentrated exactly where one
would expect if the model were responding to surface fluency rather than
substance: it is small on Steps to Reproduce (a checkable, structural judgment
that survives rewriting) and severe on Observed and Expected Behavior (free-text
sufficiency judgments that the ImproBR rewrite makes read as complete). The
bias test confirms the direction: on Improved the model scores higher than the
human in 47 of 52 disagreements and calls 111 of 139 reports reproducible. In
effect, the AI polish that makes a report look thorough also makes the evaluator
credit it, whether or not the underlying reproducibility information is really
there.

This has a concrete practical implication. Modern triage pipelines increasingly
contain LLM-generated or LLM-rewritten text, and a natural design is to chain an
LLM improver with an LLM evaluator. Our results warn that such a chain can be
self-reinforcing: the evaluator systematically over-credits the improver's
output. An LLM evaluator should therefore be validated on the specific kind of
text it will score in deployment---validating it only on raw human reports would
have overstated its reliability on the improved reports by nearly $0.25$ kappa.

\subsection{Agreement Is Uneven Across Dimensions, Consistently}
\label{subsec:disc-dimensions}

On both forms the dimension ranking is identical: Steps to Reproduce is
strongest (Raw $\kappa = 0.693$, Improved $0.654$), Expected Behavior is
intermediate, and Observed Behavior is weakest (Raw $0.289$, Improved $0.130$).
Steps to Reproduce is closest to a binary, checkable judgment---does the report
state enough to attempt the action---while Observed Behavior sufficiency
requires judging whether a free-text description is complete enough, a more
subjective distinction. That the same ordering holds on both forms, while the
softer dimensions degrade much more under rewriting, suggests an LLM-as-evaluator
is more trustworthy for structurally checkable sub-judgments than for open-ended
sufficiency judgments, and that this is a stable property of the model rather
than an artifact of one dataset.

\subsection{Coarse Screening Survives on Both Forms}
\label{subsec:disc-binarized}

The binarized check (score $\geq 4$ as reproducible) reaches $87.1\%$ accuracy
on Raw and $86.3\%$ on Improved, in both cases noticeably stronger than the
fine-grained ordinal agreement. Most of the model's disagreement with humans
happens between adjacent ordinal categories rather than between clearly
reproducible and clearly not. Practically, GPT-4o mini is more dependable as a
coarse first-pass filter than as a fine-grained 1--5 scorer---though on the
Improved form the high binary recall ($98.9\%$) comes at the cost of precision,
because the model calls almost everything reproducible, so even the coarse
filter inherits the optimistic bias when the input has been AI-polished.

\subsection{Relation to Prior Work}
\label{subsec:disc-related}

ImproBR reports that its pipeline raises the share of Executable reports from
$28.8\%$ to $67.6\%$ on this dataset~\cite{akyol2026improbr}; our human
ground truth reproduces those exact figures, confirming that we are evaluating
on the same material. Our contribution is complementary and cautionary: rather
than improving reports, we ask whether an LLM can \emph{judge} their
reproducibility with human-level reliability, and we find that it approaches
that bar on original reports but falls well short on the very reports ImproBR
improves. Combined with the LLM-as-judge literature's caution that automated
evaluators must be validated against human labels before their agreement figures
can be trusted~\cite{li2024judges}, our result sharpens that caution: validation
must be done on the same distribution of text---including AI-generated
text---that the evaluator will face in use, because agreement measured on human
reports does not transfer to machine-rewritten ones.
